\chapter{Neural Network Components: Losses}
\label{Losses}


\begin{figure}[!ht]
    \centering
    \includegraphics[width=\textwidth]{../../images/07_Neural Network Loss Functions Explained.png
    }
\end{figure}

% Document layers, loss functions, and model APIs here.

From the previous chapters, we've established
the core components of our framework, including
tensors, operations, autodiff, and neural
network components such as layers.

So far, our framework is good at performing
mathematical computations through forward and
backward passes. But how does it know whether
its predictions are correct? More importantly,
how does it know whether training is making
progress?

Our model's trainable parameters are initialized,
and gradients can be calculated through the
computational graph. But we still need a way
to measure the difference between the model's
predictions and the expected outputs.

That's where loss functions come in.

In supervised learning, a loss function measures
the discrepancy between predicted outputs and their
corresponding targets. During training, we typically
aim to minimize this loss. For each batch, the
model produces predictions during the forward pass,
and the loss function evaluates how far those
predictions are from the actual targets.

Let's see how this works in practice.

\section{Loss Functions: Mean Squared Error (MSE)}

\[
\mathrm{MSE} = \frac{1}{N}\sum_{i=1}^{N}(Y_i-\hat{Y}_i)^2
\]

where \(N\) is the number of elements over which
the squared errors are averaged, \(Y_i\) is the
target, and \(\hat{Y}_i\) is the predicted output.

For each element, we subtract the prediction from
the target and square the resulting difference. We
then calculate the mean of these squared errors
over the selected elements.

\subsection{MSE: Differentiation}

Although our autodiff engine calculates the gradients
automatically, let's look at the derivative to
understand what happens during backpropagation.

Differentiating MSE with respect to a prediction gives:

\[
\frac{\partial \mathrm{MSE}}{\partial \hat{Y}_i}
=\frac{2}{N}(\hat{Y}_i-Y_i)
\]

This gradient measures how the loss changes with
respect to each predicted value. Our autodiff
engine propagates it through the computational
graph to calculate gradients for the model's
trainable parameters.

\subsection{MSE: Putting things together}

Okey, now as we've sorted out the math part, we
can now focus on the implementation. If we think
deeply losses share common traits from their
bigger sibling layers. They have few differences

\begin{itemize}
    \item 1. Don't have trainable parameter
    \item 2. Outputs a scaler value rather than
    tensor.
\end{itemize}

So, we can follow same structure like Layers

\begin{lstlisting}[
    caption={Creating loss class},
    label={lst:layer abstruct class}
]
// file loss.hpp

class Loss {
public:
    virtual ~Loss() = default;

    virtual float calculate(std::vector<tf::tensor> inputs) = 0;
    virtual void forward(std::vector<tf::tensor> inputs) = 0;
    virtual void backward() = 0;
};
\end{lstlisting}

Although the loss function does not need to
implement differentiation manually, its backward
interface can provide access to the gradients
recorded by the computational graph. This is
useful for debugging and inspecting the
intermediate gradients.

Let's see how MeanSquaredError(MSE) is layed out.

\begin{lstlisting}[
    caption={Creating MSE class},
    label={lst:MSE class}
]
// file mse.cpp
#include <mse.hpp>

    void MSE::forward(std::vector<tf::tensor> inputs) {
      if (inputs.size() == 1) {
        training_inputs.push_back(inputs[0]);
        ...
        // initialize the intermediate tensors
        ...

        differences = target_ouput.sub(training_inputs[0]);
        loss_tensor = differences.pow(2).scale(1.0/batch_size);
      }
    };
    void MSE::backward(){
        if (!this->isGradRecorded) {
            this->isGradRecorded = true;
            Graph *g = GraphManager::instance().getCurrentGraph();
            grad_training_inputs = g,getGradTensor(trining_inputs[0]);
        }
    };
\end{lstlisting}

but sadly I've don't have any standalone
configuration to showcase how a user can
use losses. Most of its usage are inside
model. As of now simple eager execution is
not ready yet.

\section{Loss Function: Categorical Cross Entropy (CCE)}

While MSE is commonly used for regression tasks, Categorical Cross Entropy (CCE) is commonly used for classification tasks. It measures how well the predicted class probabilities match the target distribution.

For one-hot encoded targets, CCE is defined as:

\[
\mathrm{CCE}=-\frac{1}{N}
\sum_{i=1}^{N}\sum_{c=1}^{C}
Y_{ic}\log(\hat{Y}_{ic})
\]

where \(N\) is the batch size and \(C\) is the number of classes.

For each sample, we multiply the target value
for each class by the logarithm of its predicted
probability, sum over the classes, and average
over the batch. The negative sign ensures that
minimizing the loss encourages the model to
assign higher probabilities to the correct classes.

\subsection{CCE: Differentiation}

\[
\frac{\partial{CCE_\text{batch}}}{\partial{\hat(Y)}}= - \frac{1}{n}\frac{Y}{\hat{Y}}
\]

\subsection{CCE: Putting things together}

\begin{lstlisting}[
    caption={Creating CCE class},
    label={lst:CCE class}
]
// file cce.cpp
#include <cce.hpp>

    void CCE::forward(std::vector<tf::tensor> inputs) {
      if (inputs.size() == 1) {
        training_inputs.push_back(inputs[0]);
        ...
        // initialize the intermediate tensors
        ...

        log_differences = target_inputs[0].mul(training_inputs[0].log());
        loss_tensor = log_differences.scale(1.0/batch_size);
      }
    };
    void CCE::backward(){
        if (!this->isGradRecorded) {
            this->isGradRecorded = true;
            Graph *g = GraphManager::instance().getCurrentGraph();
            grad_training_inputs = g,getGradTensor(trining_inputs[0]);
        }
    };
\end{lstlisting}

\subsection{CCE: In practice}

Like MSE I've don't have eager mode calculate()
is still not ready. So I don't have much ex.
to showcase.

% Document layers, loss functions, and model APIs here.
